% ЗГЕНЕРОВАНО scripts/build_format_v2_chapter.py — не редагувати вручну.
% Джерело: editorial/cards/4.*.json, format-v2/rozdil-4.json, paraleli.json.

\naukchapter{4}{6}{Глава 4. Хто пізнав — не боїться}
\naukintro{Четверта глава найкоротша в першій половині книги: шість віршів. Вона не додає нового вчення, а перевіряє попереднє питанням про міру: чим відрізняється той, хто пізнав? Відповідь будується на порівняннях — небо й дим, чотири роди істот, — і жодного разу не на обіцянці нагороди. Останній вірш замикає главу тим, що знімає страх: не як наслідок мужності, а як наслідок знання.}
\vsec{4.1}
\vdeva{हन्तात्मज्ञस्य धीरस्य खेलतो भोगलीलया । न हि संसारवाहीकैर् मूढैः सह समानता ॥}
\vtranslit{хантатмаджняся дгірася кхелато бгоґалілая | на хі самсаравахікайр мудгайх саха саманата ||}
\vperek{О! Атмазнавцеві, мудрому, що грається грою досвідів, немає рівні з дурнями, які тягнуть сансару.}
\vcom{Построкове зіставлення: редакторська чернетка, не рецензовано.}
\vcom{Рядок а}
\vdeva{हन्तात्मज्ञस्य धीरस्य खेलतो भोगलीलया}
\vtranslit{хантатмаджняся дгірася кхелато бгоґалілая}
\vperek{О! Атмазнавцеві, мудрому, що грається грою досвідів,}
\vcom{Вигук hanta відкриває подив. Форми ātmajñasya й dhīrasya означують того, хто знає Атман; khelato та bhogalīlayā вводять образ гри досвідів. Це перший смисловий фрагмент одного речення, а не самостійна теза.}
\vcom{Рядок б}
\vdeva{न हि संसारवाहीकैर् मूढैः सह समानता}
\vtranslit{на хі самсаравахікайр мудгайх саха саманата}
\vperek{немає рівні з дурнями, які тягнуть сансару.}
\vcom{Заперечення na hi ... samānatā завершує думку про відсутність рівності; saṃsāravāhīkair mūḍhaiḥ saha називає тих, із ким здійснено порівняння. Образ тягаря сансари взято з чинної картки; атрибуція мовця між цифровими свідками лишається спірною.}
\vcom{Слово vāhīka («погонич, той, хто везе») вибране точно: сансара тут не стан, у якому перебувають, а вантаж, який тягнуть. Дурість полягає не в незнанні, а в тому, що людина впряглася. Протиставлення будується на двох способах дії: khelataḥ («того, хто грається») і vāhīkaiḥ («тих, хто тягне»). Те саме діяння, різна якість. līlā («гра») — слово з релігійної мови, де гра не протистоїть серйозності, а означає дію без потреби. Вигук hanta на початку задає тон подиву, а не докору.}
\vpar{Українська паралель (наша інтерпретація)}{Леся Українка, «Лісова пісня» (1911).}{Наша паралель: у Лесі Українки гра теж не протистоїть досвіду, а є способом жити серед нього, не застигаючи. Наше порівняння — типологічне, не твердження про джерело.}
\vfig{ilustratsii/4-1-motion.pdf}{Дві дороги: гра досвідів і тягар сансари}{Схема наша, не історична ілюстрація.}

\vsec{4.2}
\vdeva{यत् पदं प्रेप्सवो दीनाः शक्राद्याः सर्वदेवताः । अहो तत्र स्थितो योगी न हर्षम् उपगच्छति ॥}
\vtranslit{ят падам препсаво дінах шакрадях сарвадеватах | ахо татра стхіто йоґі на харшам упаґаччаті ||}
\vperek{О диво: те, до чого прагнуть убогі всі боги, починаючи з Шакри, — йог, що вже стоїть у цьому, не відчуває захоплення.}
\vcom{Построкове зіставлення: редакторська чернетка, не рецензовано.}
\vcom{Рядок а}
\vdeva{यत् पदं प्रेप्सवो दीनाः शक्राद्याः सर्वदेवताः}
\vtranslit{ят падам препсаво дінах шакрадях сарвадеватах}
\vperek{О диво: те, до чого прагнуть убогі всі боги, починаючи з Шакри, —}
\vcom{Форми yat padaṃ і prepsavaḥ описують стан, до якого прагнуть; dīnāḥ застосовано навіть до śakrādyāḥ sarvadevatāḥ. Український вигук «О диво» перенесено на початок усього речення з наступного санскритського піврядка — це переставлення для плавного читання, не дослівний порядок.}
\vcom{Рядок б}
\vdeva{अहो तत्र स्थितो योगी न हर्षम् उपगच्छति}
\vtranslit{ахо татра стхіто йоґі на харшам упаґаччаті}
\vperek{йог, що вже стоїть у цьому, не відчуває захоплення.}
\vcom{У другому піврядку aho виражає подив, tatra sthito yogī називає йога, який уже перебуває в цьому стані, na harṣam upagacchati заперечує захоплення. Порівняння прагнення та незворушності — змістова структура вірша, а не оцінка реальних релігійних громад.}
\vcom{Слово dīnāḥ («убогі») вжите про богів навмисно: убогість тут не в майні, а в прагненні — хто прагне, той має брак. Дезидератив prepsavaḥ («ті, що хочуть досягти») підкреслює саме бажання досягти, а не досягнення. Йог описаний дієприкметником sthitaḥ («той, хто стоїть») — це стан, а не подія; і саме тому harṣa («захоплення») не виникає: захоплення потребує новизни, а тут нічого нового не сталося. Вигук aho той самий, що в серії «див» глави 3, і тут він позначає межу між прагненням і перебуванням.}
\vpar{Українська паралель (наша інтерпретація)}{Григорій Сковорода, «Вдячний Еродій» (трактат).}{Наша паралель: у Сковороди прагнення саме по собі теж викриває брак, і тому досягнуте не викликає захоплення. Наша інтерпретація за будовою думки, а не вказівка на вплив.}
\vfig{ilustratsii/4-2-balance.pdf}{Дві шальки: прагнення богів і спокій того, хто вже стоїть}{Схема наша, не історична ілюстрація.}

\vsec{4.3}
\vdeva{तज्ज्ञस्य पुण्यपापाभ्यां स्पर्शो ह्य् अन्तर् न जायते । न ह्य् आकाशस्य धूमेन दृश्यमाना ऽपि संगतिः ॥}
\vtranslit{таджджняся пуняпапабгям спаршо хй антар на джаяте | на хй акашася дгумена дрішямана ’пі самґатіх ||}
\vperek{Того, хто осягнув Істину, не зачіпають ані заслуги, ані провини. Так само й небо не зливається з димом, хоча на око вони ніби поєднані.}
\vcom{Построкове зіставлення: редакторська чернетка, не рецензовано.}
\vcom{Рядок а}
\vdeva{तज्ज्ञस्य पुण्यपापाभ्यां स्पर्शो ह्य् अन्तर् न जायते}
\vtranslit{таджджняся пуняпапабгям спаршо хй антар на джаяте}
\vperek{Того, хто осягнув Істину, не зачіпають ані заслуги, ані провини.}
\vcom{Родовий tajjñasya («того, хто знає Те») й подвійна форма puṇyapāpābhyāṃ у конструкції sparśo hy antar na jāyate пояснюють заперечений дотик заслуг і провин. «Істина» — чинний перекладацький відповідник, а не доведене буквальне значення займенника tat поза контекстом.}
\vcom{Рядок б}
\vdeva{न ह्य् आकाशस्य धूमेन दृश्यमाना ऽपि संगतिः}
\vtranslit{на хй акашася дгумена дрішямана ’пі самґатіх}
\vperek{Так само й небо не зливається з димом, хоча на око вони ніби поєднані.}
\vcom{Порівняння na hy ākāśasya dhūmena та dṛśyamānā 'pi saṃgatiḥ розрізняє видиму суміжність диму й неба та їхню сутнісну єдність. Це аналогія самого вірша; вона не є твердженням про фізичну природу атмосфери.}
\vcom{У вірші 4.3 використано порівняння неба і диму, щоби розрізнити видимий зв'язок та сутнісне ототожнення. Зовнішнє спостереження диму на тлі неба ще не означає, що небо набуває властивостей диму. Аналогічно фігура «того, хто знає Те» в недвоїстому вченні не визначається заслугами чи провинами як власною сутністю. З цього не випливає юридичний або етичний дозвіл чинити шкоду: вірш формулює філософську тезу про природу свідка, а не правила поведінки. У робочому перекладі «осягнув Істину» — тлумачна експлікація займенника tat, можлива буквальніша редакція — «той, хто знає Те»; на рецензії слід обрати один варіант. Потрібно окремо звірити різницю «з'єднання» / «дотик» і мовну форму фінального іменника saṅgati / saṃgatiḥ. **Атрибуція мовця не усталена за звіреними електронними свідками:** GRETIL/TextGrid перед 4.1 подає «aṣṭāvakra uvāca», а SARIT в заголовку глави 4 — «janaka uvāca». Поки не залучено незалежне друковане видання чи рукописні свідчення, не слід беззастережно приписувати твердження Аштавакрі або Джанаці. Два електронні свідки не є критичним виданням; коментарі сучасних авторів тут не передруковано.}
\vpar{Українська паралель (наша інтерпретація)}{Памфіл Юркевич, «Серце і його значення в духовному житті людини» (1860).}{Наша паралель: коли вірш каже, що небо не зливається з димом, він відділяє те, що проходить, від того, на чому воно проходить. Юркевич розводить серце й привхідне в людині так само рішуче. Наше порівняння — за способом розрізняти.}
\vfig{ilustratsii/4-3-space.pdf}{Небо й дим: видно разом, а не зливаються}{Схема наша, не історична ілюстрація.}
\vfig{ilustratsii-naratyvni/4-3-dym-i-nebo.pdf}{Дим від вогнища проходить крізь небо, але не змішується з ним}{Наша ілюстрація до вірша 4.3 — не історичне зображення.}

\vsec{4.4}
\vdeva{आत्मैवेदं जगत् सर्वं ज्ञातं येन महात्मना । यदृच्छया वर्तमानं तं निषेद्धुं क्षमेत कः ॥}
\vtranslit{атмайведам джаґат сарвам джнятам єна махатмана | ядріччая вартаманам там нішеддгум кшамета ках ||}
\vperek{Хто здатен стримати того, хто пізнав увесь світ як єдиний Атман і живе так, як складаються обставини?}
\vcom{Построкове зіставлення: редакторська чернетка, не рецензовано.}
\vcom{Рядок а}
\vdeva{आत्मैवेदं जगत् सर्वं ज्ञातं येन महात्मना}
\vtranslit{атмайведам джаґат сарвам джнятам єна махатмана}
\vperek{Хто здатен стримати того, хто пізнав увесь світ як єдиний Атман}
\vcom{Конструкція ātmaivedaṃ jagat sarvaṃ jñātaṃ yena mahātmanā називає пізнання всього світу як Атмана. Чинний український переклад починається з питання «Хто здатен стримати...», яке в санскриті завершальне: український поділ передає смислові половини, але не дзеркальний порядок слів.}
\vcom{Рядок б}
\vdeva{यदृच्छया वर्तमानं तं निषेद्धुं क्षमेत कः}
\vtranslit{ядріччая вартаманам там нішеддгум кшамета ках}
\vperek{і живе так, як складаються обставини?}
\vcom{yadṛcchayā vartamānaṃ описує життя відповідно до обставин, а taṃ niṣeddhuṃ kṣameta kaḥ формулює риторичне питання, хто може стримати цю людину. Тут не виводиться дозвіл нехтувати правами інших чи відповідальністю.}
\vcom{Вірш поєднує дві теми: пізнання недвоїстої природи світу та свободу від нав'язаної поведінкової схеми. Питання «хто здатен стримати?» є риторичним і характеризує свободу того, хто усвідомив світ як Атман. Слово yadṛcchayā не слід автоматично прирівнювати до примхливого бажання або безвідповідальності: у цьому контексті воно допускає читання «за збігом обставин, природним чином». Українське «єдиний Атман» — смислове відтворення підсилювального eva, а не окремий прикметник оригіналу; перед остаточним виданням перевірити, чи краще писати «сам Атман». Важливо, що vartamānaṃ taṃ граматично підтримує тлумачення про людину, яка живе спонтанно, а не тезу, ніби «світ робить усе, що йому хочеться». Цей коментар — редакторська інтерпретація, не незалежна наукова рецензія. Додатковий текстологічний блокер: GRETIL і SARIT не збігаються в атрибуції мовця всієї четвертої глави, тому цей вірш не приписано беззастережно жодному співрозмовникові.}
\vpar{Українська паралель (наша інтерпретація)}{Михайло Коцюбинський, «Intermezzo» (1908).}{Наша паралель: у Коцюбинського той, хто бачить ціле, теж не піддається керуванню ззовні. Наше порівняння — типологічне.}
\vfig{ilustratsii/4-4-witness.pdf}{Центр, якого не втримати: він бачить усе як єдине}{Схема наша, не історична ілюстрація.}

\vsec{4.5}
\vdeva{आब्रह्मस्तम्बपर्यन्ते भूतग्रामे चतुर्विधे । विज्ञस्यैव हि सामर्थ्यं इच्छा ऽनिच्छाविवर्जने ॥}
\vtranslit{абрахмастамбапар'янте бгутаґраме чатурвідге | віджнясяйва хі самартхям ічча ’ніччавіварджане ||}
\vperek{Серед чотирьох родів істот — від Брахми до стеблини трави — лише той, хто пізнав, спроможний полишити і бажання, і небажання.}
\vcom{Построкове зіставлення: редакторська чернетка, не рецензовано.}
\vcom{Рядок а}
\vdeva{आब्रह्मस्तम्बपर्यन्ते भूतग्रामे चतुर्विधे}
\vtranslit{абрахмастамбапар'янте бгутаґраме чатурвідге}
\vperek{Серед чотирьох родів істот — від Брахми до стеблини трави —}
\vcom{ābrahmastambaparyante bhūtagrāme caturvidhe задає межі сукупності істот та згадує чотири роди. Образ «від Брахми до стеблини» передає контраст міфологічного найвищого й звичайної рослини; докладна класифікація потребує окремого текстологічного підтвердження.}
\vcom{Рядок б}
\vdeva{विज्ञस्यैव हि सामर्थ्यं इच्छा ऽनिच्छाविवर्जने}
\vtranslit{віджнясяйва хі самартхям ічча ’ніччавіварджане}
\vperek{лише той, хто пізнав, спроможний полишити і бажання, і небажання.}
\vcom{vijñasyaiva hi sāmarthyaṃ пов'язує спроможність із тим, хто знає; icchā 'nicchāvivarjane позначає полишення бажання й небажання. Український «лише» зберігає eva, але не проголошує знання соціальним привілеєм.}
\vcom{Вірш протиставляє всеосяжний образ світу істот — від найвищого міфологічного персонажа до найскромнішої рослини — стану знання. Ключовий вислів vijñasyaiva hi sāmarthyam указує не на соціальний ранг, а на спроможність «того, хто знає». Пара icchā / anicchā охоплює внутрішні схильності до прийняття й відкидання; переклад «бажання й небажання» ближчий до буквального заперечного словотвору, тоді як «бажання й відраза» є сильнішою філософською інтерпретацією, яку варто перевірити в незалежному коментарі. «Чотири роди» не слід подавати як сучасну біологічну систематику: вірш використовує традиційну космологічну категоризацію, але не перелічує її членів. У GRETIL/TextGrid перед 4.1 позначено «aṣṭāvakra uvāca», тоді як у SARIT заголовок четвертої глави має «janaka uvāca». Через невирішену суперечність заголовків коментар не приписує ці слова певному мовцеві. Це власне робоче пояснення за двома електронними транскрипціями, не критичне видання й не незалежна фахова рецензія.}
\vpar{Українська паралель (наша інтерпретація)}{Іван Франко, «Мойсей» (1905).}{Наша паралель: у Франка вибір, що не зводиться до бажання й небажання, теж постає як те, на що здатний не кожен. Наше порівняння — за будовою думки, не про вплив.}
\vfig{ilustratsii/4-5-neti.pdf}{Бажання й небажання відкинуто разом}{Схема наша, не історична ілюстрація.}

\vsec{4.6}
\vdeva{आत्मानम् अद्वयं कश्चिज् जानाति जगद् ईश्वरं । यद् वेत्ति तत् स कुरुते न भयं तस्य कुत्रचित् ॥}
\vtranslit{атманам адваям кашчідж джанаті джаґад ішварам | яд ветті тат са куруте на бгаям тася кутрачіт ||}
\vperek{Хтось знає Атман — недвоїстий, володаря світу. Що він знає, те й робить, — і страху в нього немає ніде.}
\vcom{Построкове зіставлення: редакторська чернетка, не рецензовано.}
\vcom{Рядок а}
\vdeva{आत्मानम् अद्वयं कश्चिज् जानाति जगद् ईश्वरं}
\vtranslit{атманам адваям кашчідж джанаті джаґад ішварам}
\vperek{Хтось знає Атман — недвоїстий, володаря світу.}
\vcom{ātmānam advayaṃ ... jagad īśvaraṃ є об'єктом пізнання; kaścij jānāti називає неназваного «когось». У словосполученні про володаря світу збережено чинну редакцію без привнесення нових богословських атрибуцій.}
\vcom{Рядок б}
\vdeva{यद् वेत्ति तत् स कुरुते न भयं तस्य कुत्रचित्}
\vtranslit{яд ветті тат са куруте на бгаям тася кутрачіт}
\vperek{Що він знає, те й робить, — і страху в нього немає ніде.}
\vcom{yad vetti tat sa kurute зіставляє знання і дію; na bhayaṃ tasya kutracit заперечує страх будь-де. Це формула філософського тексту, не перевірене твердження про психологічну безстрашність конкретної людини.}
\vcom{kaścit («хтось») на початку — не випадкове слово: вірш не обіцяє, що це стосується кожного, і не називає його. Друга половина дає формулу yad vetti tat kurute («що знає, те й робить»), і вона сильніша за звичне «чини за знанням»: тут знання й дія не узгоджуються, а ототожнюються. kutracit з na («ніде») закриває простір повністю — страх не має місця, а не просто відсутній тут. Це останній вірш глави, і він повертає до її початку: 4.1 протиставляв того, хто грається, тому, хто тягне; тут сказано, чому — бо немає страху, а отже немає й потреби тягнути.}
\vpar{Українська паралель (наша інтерпретація)}{Василь Стус, «Час творчості».}{Наша паралель: у Стуса теж є стан, у якому людина робить те, що знає, і не має страху ніде. Наша асоціація, а не твердження про джерело.}
\vfig{ilustratsii/4-6-witness.pdf}{Той, хто знає недвоїсте, робить те, що знає}{Схема наша, не історична ілюстрація.}

